\documentclass[12pt]{article}
\usepackage[margin=0.6in,top=1.15in,bottom=0.55in]{geometry}
\usepackage{lmodern}
\usepackage{microtype}
\usepackage{titlesec}
\usepackage{enumitem}
\usepackage[hidelinks]{hyperref}
\usepackage{../../latex/grader-worksheet}
\setlength{\parindent}{0pt}
\setlength{\parskip}{0.25em}
\titleformat{\section}{\large\bfseries\scshape}{}{0pt}{}
\GraderSetup{assignment-id=U1L18LE,template-revision=1}
\begin{document}
\begin{center}
{\LARGE\bfseries Week 18 Lit Example Worksheet}\par\vspace{0.05em}
{\large\scshape Revising a Shakespeare Judgment}\par\vspace{0.12em}
\rule{0.78\textwidth}{0.5pt}
\end{center}
\textbf{Purpose:} Use the reader to build a visible revision trail from an early Shakespeare judgment to a precise, evidence-based logical analysis.

\section*{Part 1: Macbeth --- Read the Model}
\textbf{Question 1.1}\\[-0.15em]
In the assisted model, what does the early claim notice, and what does it oversimplify about the apparition's words?
\GraderAnswerBox{Part 1 Q1}{2.2in}

\textbf{Question 1.2}\\[-0.15em]
State the model's revised judgment and explain why it is stronger than ``the witches lie.''
\GraderAnswerBox{Part 1 Q2}{2.35in}

\newpage
\section*{Part 2: Macbeth --- Name the Revision Trail}
\textbf{Question 2.1}\\[-0.15em]
Identify the evidence, inference, and action in Macbeth's response to the promise. Which logical habit best tests the inference?
\GraderAnswerBox{Part 2 Q1}{2.45in}

\textbf{Question 2.2}\\[-0.15em]
Write a reflective caption for the model revision: original, diagnosis, substantive change, and what the change shows about judgment.
\GraderAnswerBox{Part 2 Q2}{2.55in}

\newpage
\section*{Part 3: Brutus --- Diagnose and Revise}
\textbf{Question 3.1}\\[-0.15em]
An early response says, ``Brutus cares about Rome, so his decision is logical.'' Diagnose one exact weakness in that response using the passage.
\GraderAnswerBox{Part 3 Q1}{2.4in}

\textbf{Question 3.2}\\[-0.15em]
Write a revised judgment that preserves what the passage supports while testing the move from possible danger to necessary death.
\GraderAnswerBox{Part 3 Q2}{2.55in}

\newpage
\section*{Part 4: Caption and Portfolio Connection}
\textbf{Question 4.1}\\[-0.15em]
Write a reflective caption for the Brutus revision. Name the original limit, the Whately criterion, the textual change, and the stronger judgment now possible.
\GraderAnswerBox{Part 4 Q1}{2.75in}

\textbf{Question 4.2}\\[-0.15em]
Compare the two revision trails. What recurring habit of disciplined reading do both make visible, and how could that habit improve one Shakespeare artifact in your portfolio?
\GraderAnswerBox{Part 4 Q2}{2.5in}
\end{document}
